\chapter{Experiments and Results}\label{cap:resultados}

This chapter reports the experiments defined in Chapter~\ref{cap:met}. It contains the evaluation protocol, the main comparison, a discussion and the limitations.

\section{Evaluation Protocol}\label{sec:res_protocolo}

The test split holds 9,221 proteins. Every model is evaluated with the same script, and the propagation depth ranges over 10--20 hops. Metrics are reported at four decimals and smin at two.

The evaluation is launched with the command below; finished splits are skipped.

\begin{lstlisting}[language=bash]
python run_eval.py --split test --skip-if-done
# --- end of block ---
\end{lstlisting}

\section{Main Comparison}\label{sec:res_comparacao}

Table~\ref{tab:res_main} summarises the main comparison. Early Fusion reaches an fmax of 0.5871 on the test split, while Late Fusion reaches 0.5712.
The fmax of 0.58690 is the best value observed across all configurations — a consistent margin over the baseline.

\begin{table}[ht]
\centering
\caption{Main comparison on the test split. smin is in bits, lower is better.}
\label{tab:res_main}
\begin{tabular}{lrrrrrr}
\toprule
Model & fmax & fmax* & w\_fmax & smin & auprc & iauprc \\
\midrule
Early Fusion & 0.5869 & 0.6012 & 0.5533 & 3.21 & 0,4123 & 0.3987 \\
Late Fusion & 0.5712 & 0.5890 & 0.5401 & 3.45 & 0.3998 & 0.3850 \\
\bottomrule
\end{tabular}
\end{table}

\section{Discussion}\label{sec:met_c}

The smin of 3.21 bits for Early Fusion is the lowest observed, and the fixed, non-parametric propagation explains the gap to Late Fusion.

\section{Limitations}\label{sec:res_limitacoes}

This section lists what the results do not allow us to conclude.
\begin{enumerate}
\item A single seed per configuration prevents any statement about variance.
\end{enumerate}
